\scanpage{88}
On trouve donc un diagramme de suites exactes
\[
\diagram{\begin{tikzcd}[ampersand replacement=\&,column sep=large,row sep=large]
 \& \& \& 1 \arrow[d] \& 1 \arrow[d] \& \\
 \& \& \& \deleted{I_{\unsure{s'}}}\inserted{J} \arrow[r,phantom,"\simeq"] \arrow[d] \& I_{s'}\simeq T \arrow[d] \& \\
 \& 1 \arrow[r] \& \pi \arrow[r] \arrow[d,equal] \& E_{U_{\widetilde K}} \arrow[r] \arrow[d] \& \underset{\substack{\Vert\\\operatorname{Gal}(\overline K/\widetilde K)}}{E_{\widetilde K}} \arrow[r] \arrow[d] \& 1 \\
(26) \& 1 \arrow[r] \& \pi \arrow[r] \& E_{U_{\unsure{\widetilde O|\widetilde S}}} \arrow[r] \arrow[d] \& \underset{\substack{\Vert\\\operatorname{Gal}(\overline k/k)}}{E_{\unsure{\widetilde O|\widetilde S}}} \arrow[r] \arrow[d] \& 1 \\
 \& \& \& 1 \& 1 \&
\end{tikzcd}}
\]
le noyau \inserted{$J$} de $E_{U_{\widetilde K}}-E_{U_{\unsure{\widetilde O|\widetilde S}}}$ s'envoyant isomorphiquement par $E_{U_{\widetilde K}}\longrightarrow E_{\widetilde K}$ sur le sous-groupe d'inertie $I_{s'}$, lui-$\unsure{\overline m|\overset{m}{=}}$ can. isom. à $T$.

La section de \inserted{$E_{U_{\widetilde K}}$ sur} $E_{\widetilde K}\simeq N_{s'}$, induite par la section de départ de $E_{U_\eta}$ sur $E_K$ (elle $\unsure{\overline m|\overset{m}{=}}$ définie par \inserted{le pt $K$-rat.} $f_K$ de $U_\eta$ sur $\eta$) n'est autre bien sûr que celle définie par le pt $\widetilde K$-rationnel correspondant $f_{\widetilde K}$ de $U_{\widetilde K}$. Son composé avec $E_{U_{\widetilde K}}-E_{U_{\unsure{\widetilde O|\widetilde S}}}$ envoie $I\simeq T$ dans $\pi$, on a vu comment : essentiellement par un $\kappa_i\circ(\mu.\operatorname{id}_T)$, où $\kappa_i:T-\pi$ est un hom.-local, relatif à l'indice $i$.

Ainsi \inserted{le groupe de décomposition} $E_{\widetilde K}\deleted{\unsure{(s')}}$ s'envoie dans le normalisateur de $L_i=\kappa_i(T)$, dans $E_{U_{\unsure{\widetilde O|\widetilde S}}}$. D'ailleurs
\margin{87}
